\section{Related Work}
\label{sec:related}

\subsection{Fact Verification in Low-Resource Languages}
Fact verification benchmarks have expanded substantially since the foundational FEVER challenge~\citep{thorne2018fever}, which formalized verification as evidence retrieval over Wikipedia followed by natural language inference (NLI) classification into \textsc{Supported}, \textsc{Refuted}, or \textsc{Not Enough Info} (NEI). Subsequent benchmarks incorporated multi-hop tabular reasoning~\citep{aly2021feverous}, real-world claim styles~\citep{schuster2019fever}, and causal structure modeling~\citep{si-etal-2024-checkwhy}. Recently, retrieval-augmented verification has evolved toward synthesizing contrastive arguments~\citep{yue-etal-2024-retrieval} and deploying specialized, lightweight verifiers such as MiniCheck~\citep{tang-etal-2024-minicheck} to ground LLM outputs. Crucially, recent findings establish that evidence retrieval remains the primary empirical bottleneck for factual verification~\citep{zheng-etal-2024-evidence}. While cross-lingual benchmarks such as TyDi QA~\citep{clark2020tydi} have broadened language coverage, Central Asian languages remain critically underserved. Prior non-Indo-European datasets almost universally overlook agglutinative morphology, allowing lexical overlap heuristics to dominate baseline model performance. Kazakh-FEVER 3K resolves this gap by coupling gold evidence alignments from Kazakh encyclopedic corpora with a morphologically-aware hybrid retrieval and NLI pipeline.

\subsection{Morphologically-Aware Turkic NLP}
The Turkic language family, comprising over 35 distinct languages, is characterized by extensive agglutinative affixation and strict vowel harmony. Early computational work on Kazakh centered on rule-based finite-state transducers (FSTs) and morphological analyzers~\citep{makhambetov2013kazakh,bekmanova2021kazakh}. With the ascent of large-scale language models, multilingual transformers such as mBERT~\citep{devlin2019bert} and XLM-RoBERTa~\citep{conneau2020xlmr}, as well as monolingual models like Kaz-RoBERTa~\citep{sagyndyk2025kazroberta}, have advanced Kazakh representation learning. Recent initiatives have further expanded Kazakh resources through culturally grounded instruction tuning~\citep{laiyk-etal-2025-instruction} and sentiment benchmarks~\citep{yeshpanov2024kazsandra}. However, standard subword tokenizers frequently segment Kazakh inflectional chains into semantically arbitrary fragments, severely obscuring verbal negation and evidential markers. Recent work confirms that incorporating explicit morphological segmentation and morpheme boundary inductive biases is vital to mitigate data sparsity and stem distortion in agglutinative NLP~\citep{guo2026kazakh}. Our architecture unifies morpheme-aware sparse retrieval with dense representations to overcome these tokenizer pathologies.

\subsection{AI-Generated Text Detection and Trust Modeling}
The rapid proliferation of high-capacity generative language models~\citep{touvron2023llama} has intensified the threat of plausible synthetic disinformation. Existing detection paradigms primarily leverage zero-shot probability curvature, such as DetectGPT~\citep{mitchell2023detectgpt} and Fast-DetectGPT~\citep{bao2024fastdetectgpt}, or scoring perplexity ratios across paired models via Binoculars~\citep{hans2024binoculars}. Concurrently, comprehensive benchmarks like RAID~\citep{wang2024raid} reveal severe performance degradations across detectors facing adversarial perturbation or out-of-domain distribution shifts. Crucially, existing methods treat synthetic text detection and fact verification as isolated silos. Yet, as demonstrated by \citet{guan-etal-2024-language}, language models frequently hallucinate while simultaneously excelling at fact verification tasks, establishing that stylistic machine generation ($T_{\text{gen}}$) and epistemic truth ($T_{\text{fact}}$) are mathematically orthogonal dimensions. In real-world environments, human text can convey malicious falsehoods, while synthetic generation can articulate grounded truths. Our continuous 2D Trust Matrix resolves this decoupling by projecting veracity and authenticity predictions into a unified epistemic coordinate system.
